\section{Roots of unity}\label{sec:roots}

Let $\zeta_n=e^{2\pi i/n}$, let $\mu_n$ be the group of $n$th roots of unity, and put
$S_n=S(\{0\}\cup\mu_n)$.  The case $n=6$ is the alphabet $\{0,\pm1,\pm\omega,\pm\omega^2\}$ with
$\omega=e^{2\pi i/3}$, which contains the 33-ray set that Cabello calls the simplest KS
set~\cite{Cabello2025simplest}.

\begin{lemma}\label{lem:sums}
A sum of two elements of $\mu_n$ can vanish if and only if $2\mid n$.  A sum of three can vanish
if and only if $3\mid n$, and then the three terms are $\beta,\beta\omega,\beta\omega^2$ for some
$\beta\in\mu_n$.
\end{lemma}

\begin{proof}
$\zeta^a+\zeta^b=0$ iff $\zeta^{b-a}=-1$, and $-1\in\mu_n$ iff $2\mid n$.  For three terms divide
by the first: $1+e^{i\theta}+e^{i\phi}=0$.  The imaginary part gives $\phi=-\theta$ or
$\phi=\pi+\theta$; the second makes the real part equal to~1.  With $\phi=-\theta$ the real part
gives $\cos\theta=-\tfrac12$, so the other two terms are $\omega$ and $\omega^2$, which lie in
$\mu_n$ iff $3\mid n$.
\end{proof}

Hence two rays of $S_n$ are orthogonal only if the number of coordinates in which both are
nonzero is~0, or~2 with $2\mid n$, or~3 with $3\mid n$.

\begin{theorem}\label{thm:6n}
$S_n$ is KS-uncolorable if and only if $6\mid n$.
\end{theorem}

\begin{proof}
\emph{If.}  For $6\mid n$ we have $\mu_6\subseteq\mu_n$ and so $S_6\subseteq S_n$.  The pool
$S_6$ is uncolorable (Table~\ref{tab:cyc}; it contains the 33-ray set
of~\cite{Cabello2025simplest}).  Apply Lemma~\ref{lem:restrict}.

\emph{Only if.}  Let $6\nmid n$ and write $e_1,e_2,e_3$ for the axis rays.

\emph{Case 1: $2\nmid n$ and $3\nmid n$.}  Orthogonal rays have disjoint supports, so the only
triad is $\{e_1,e_2,e_3\}$.  Give $e_1$ the value~1 and every other ray~0.

\emph{Case 2: $2\mid n$, $3\nmid n$.}  Two rays with three common nonzero coordinates are never
orthogonal.  No triad contains an all-nonzero ray~$a$: each of its two partners would share
exactly two nonzero coordinates with~$a$ and so have support of size~2; two such supports meet,
so the partners, being orthogonal to each other, would have the same support; and then $a$ would
be orthogonal to a basis of that coordinate plane and vanish on it.  Rays whose supports are
different two-element sets share one coordinate and are not orthogonal.  So every triad is
$\{e_1,e_2,e_3\}$ or $\{e_k,v,v'\}$ with $v,v'$ in the coordinate plane orthogonal to~$e_k$.  In
such a plane the rays are $(1,\zeta^p)$, and $(1,\zeta^p)\perp(1,\zeta^q)$ iff $q-p\equiv n/2$: a
perfect matching.  Give $e_1$ the value~1 and $e_2,e_3$ the value~0; give~0 to every ray of the
plane orthogonal to $e_1$; in each of the other two coordinate planes give~1 to one ray of each
matched pair and~0 to the other; give~0 to every all-nonzero ray.  Rays in different coordinate
planes are not orthogonal, and no ray with value~1 in those two planes is orthogonal to~$e_1$, so
no two rays with value~1 are orthogonal, and each triad has exactly one~1.

\emph{Case 3: $3\mid n$, $2\nmid n$.}  Orthogonality needs common support of size 0 or~3.  A ray
with exactly one zero coordinate is therefore orthogonal only to the axis ray on that coordinate,
and lies in no triad.  For an all-nonzero ray normalise $v=(1,a,b)$.  By Lemma~\ref{lem:sums} an
all-nonzero $w$ is orthogonal to $v$ iff $w=w_{\pm}=(1,a\omega^{\pm1},b\omega^{\mp1})$, and
$\braket{w_+}{w_-}=1+\omega^{-2}+\omega^{2}=0$.  So the all-nonzero rays fall into disjoint
triads, with no other orthogonalities among them and none to a ray with a zero coordinate.  Give
the value~1 to one ray of each such triad and to $e_1$, and~0 to everything else.
\end{proof}

For $n\ge3$ the theorem is also proved in Lean~4 with the mathlib library, in the file
\texttt{lean/Cyclotomic6n.lean} of the repository, with no unproved steps and no axioms
beyond the standard three.  There the uncolorability of $S_6$ is checked rather than cited: 33 rays of $S_6$ are
written over $\Z[\omega]$, their orthogonalities are recomputed in that ring, and a case-split
refutation of every coloring is checked by the Lean kernel.  For $n=1,2$ the pool lies in
$S(\{0,\pm1\})$, which is colorable.

The two primes do different work.  Divisibility by~3 creates triads: at $n=9$ and $n=15$ the pool
has 28 and 76 of them and is colorable because they are pairwise disjoint.  Divisibility by~2
creates the two-term cancellations that couple triads.  Only together do the constraints
interlock.  The fields $\Q(\zeta_3)$ and $\Q(\zeta_6)$ coincide, while $S_3$ is colorable and
$S_6$ is not: Theorem~\ref{thm:6n} is a statement about alphabets, not about fields, and it has
no formulation in terms of the Galois group of the field (the one offered in version~9 was
false).

\begin{table}[t]
\centering
\caption{Cyclotomic pools, exact arithmetic (exponents modulo $n$), colorability by SAT.  The
check was run for all $2\le n\le30$ and agrees with Theorem~\ref{thm:6n} in every case, as do the
three structural statements used in Cases 1--3; a selection is shown.}
\label{tab:cyc}
\begin{tabular}{rrrrc@{\qquad\qquad}rrrrc}
\toprule
$n$ & Rays & Pairs & Triads & Colorable & $n$ & Rays & Pairs & Triads & Colorable\\
\midrule
3 & 21 & 21 & 4 & yes & 12 & 183 & 633 & 67 & \textbf{no}\\
4 & 31 & 69 & 7 & yes & 15 & 273 & 273 & 76 & yes\\
5 & 43 & 18 & 1 & yes & 16 & 307 & 843 & 25 & yes\\
6 & 57 & 174 & 22 & \textbf{no} & 18 & 381 & 1380 & 136 & \textbf{no}\\
9 & 111 & 111 & 28 & yes & 20 & 463 & 1293 & 31 & yes\\
\bottomrule
\end{tabular}
\end{table}

\begin{fact}[Eisenstein pool]\label{fact:eisenstein}
The 57-ray pool $S_6$ has 174 orthogonal pairs and 22 triads.  Its minimum uncolorable subsets
have 33 rays, certified, and there are exactly three of them, each with 78 pairs and 14 bases.
The pool is not closed under completion: $(1,1,0)$ and $(1,\bar1,1)$ are orthogonal and their
third ray $(1,\bar1,\bar2)$ is not in $S_6$.  The completed pool has 165 rays, 390 pairs and 130
triads, and its minimum is also 33, certified.
\end{fact}

Since every $S_n$ with $6\mid n$ contains $S_6$, its minimum is at most~33.  We have not
certified the minimum for $n>6$.

\begin{remark}
Cortez, Morales and Reyes~\cite{CortezMoralesReyes2022} showed that $\Z[1/6]$ is the smallest
ring $\Z[1/N]$ over which the partial ring of symmetric $3\times3$ matrices admits no algebraic
hidden state.  The number~6 appears there
and in Theorem~\ref{thm:6n}, but the two statements concern different objects (matrix entries of
projectors in one, coordinates of rays in the other) and we know of no map from one to the
other.  Version~9 presented this as a quantitative connection; it is a coincidence until someone
proves otherwise.
\end{remark}
